\structure{ОСНОВНАЯ~ЧАСТЬ}

В основной части работы последовательно рассматриваются как теоретические,
так и практические аспекты предложенного подхода. Сначала описывается методика
верификации, основанная на сопоставлении трасс исполнения реализации с
формальной спецификацией на языке TLA+. Затем приводится архитектура
разработанной распределённой системы и формальная спецификация алгоритма
консенсуса Raft. Завершает основную часть описание процедуры проверки
соответствия между реализацией и спецификацией, включая используемые
инструменты, ход экспериментов и выявленные расхождения.

\section{Методика верификации}

Предлагаемая методика верификации основана на сопоставлении трасс исполнения
программы с формальной спецификацией, написанной на языке TLA+. Схема подхода
представлена в Приложении A.

Исходными артефактами являются программная реализация алгоритма и формальная
спецификация этого алгоритма на языке TLA+. Спецификация задаёт множество
допустимых состояний системы и переходов между ними, однако не содержит
сведений о конкретной реализации. Реализация, в свою очередь, выполняется в
распределённой среде, где возможны различные порядки доставки сообщений и
недетерминированные сценарии выполнения.

Для сопоставления поведения реализации со спецификацией используется
инструментирование кода. В реализацию внедряется библиотека трассировки,
которая фиксирует события, соответствующие переходам спецификации, а также
изменения переменных спецификации. При выполнении системы каждый процесс
(координатор и участники) формирует собственную трассу в формате JSON,
содержащую последовательность операций обновления логических переменных модели.
Полученные локальные трассы агрегируются в единый глобальный файл, отражающий
наблюдаемую последовательность изменений состояния всей распределённой системы.
Таким образом, формируется абстрактное представление исполнения, выраженное в
терминах переменных TLA+, а не в терминах низкоуровневых операций кода.

Для проверки корректности трассы используется дополнительная спецификация,
построенная на основе исходной TLA+ спецификации. Она расширяет её таким
образом, чтобы поведение модели было ограничено наблюдаемой последовательностью
изменений. Иными словами, трасса интерпретируется как частичное ограничение на
допустимые переходы модели.

Далее средство модельной проверки TLC используется для проверки существования
поведения спецификации, согласованного с данной трассой. Если TLC подтверждает
корректность, это означает, что наблюдаемое исполнение допустимо спецификацией.
В противном случае обнаруживается расхождение между реализацией и моделью, и
TLC предоставляет информацию, позволяющую локализовать проблему.

